\section{信息瓶颈与任务相关表示}
\label{sec:07-Information-Theory/06-Information-and-Learning/03}

构造一个最省事的例子。标签 \(Y\) 和一百个噪声比特 \(N\) 都是相互独立的公平比特，输入是把它们拼在一起：

\[
X=(Y,N).
\]

表示直接照抄输入，\(T=X\)，那它扛着 \(101\) 比特的输入信息，其中只有 \(1\) 比特和任务有关。换成 \(T=Y\)，表示只剩 \(1\) 比特，任务信息一点没丢。

所以``信息越多越好''不是表示学习的目标。好表示要扔掉与任务无关的变化，同时留住能预测目标的那部分。信息瓶颈把这句话写成了一个可以优化的式子——这大概是``什么叫好表示''少数几个能落到纸面上的定义之一。

\subsection{压缩与预测的两个互信息}

给定联合分布 \(p(x,y)\)，从 \(X\) 随机生成表示 \(T\)，编码器写作 \(p(t\given x)\)。这里有一条结构假设：\(T\) 只看 \(X\)，不直接接触 \(Y\)，于是三者构成 Markov 链 \(Y\to X\to T\)。

两个愿望互相拉扯。一边希望 \(I(X;T)\) 小，表示没有囤积太多输入细节；一边希望 \(I(T;Y)\) 大，表示还留着任务相关的信息。写成约束形式是在 \(I(T;Y)\ge r\) 之下最小化 \(I(X;T)\)，实际更常用拉格朗日形式：

\[
\mathcal L_{\mathrm{IB}}=I(X;T)-\beta\,I(T;Y),\qquad \beta\ge 0.
\]

\(\beta\) 小的时候压缩说了算，极端情况下 \(T\) 干脆与 \(X\) 独立；\(\beta\) 大的时候预测说了算。扫过一遍 \(\beta\)，得到的不是某个普适的最佳表示维数，而是一条前沿。

天花板由数据处理不等式给出：既然 \(Y\to X\to T\)，就有 \(I(T;Y)\le I(X;Y)\)。表示造不出输入里不存在的标签信息，它能做的是重新组织信息，让有限容量的模型更容易用上，顺带扔掉妨碍泛化的无关变化。这条不等式的证明与它在特征工程里的用法见\hyperref[sec:07-Information-Theory/02-Joint-Information/04]{``数据处理不等式与充分表示''}。

\begin{center}
\begin{tikzpicture}[x=1.15cm,y=1.25cm,font=\small]
  \fill[blue!7] (0,0) -- plot[domain=0:5,samples=70]
    (\x,{2.4*(1-(1-\x/5)*(1-\x/5))}) -- (6.0,2.4) -- (6.0,0) -- cycle;
  \draw[->,black!55] (0,0) -- (6.5,0) node[right,font=\footnotesize] {\(I(X;T)\)};
  \draw[->,black!55] (0,0) -- (0,3.15) node[above,font=\footnotesize] {\(I(T;Y)\)};
  \draw[black!35,dashed] (0,2.4) -- (6.3,2.4);
  \node[left,font=\footnotesize,black!70] at (-0.06,2.4) {\(I(X;Y)\)};
  \draw[black!35,dashed] (0,0) -- (2.4,2.4);
  \node[right,font=\footnotesize,black!55,rotate=45] at (1.05,1.12) {\(I(T;Y)\le I(X;T)\)};
  \draw[blue!75,very thick] plot[domain=0:5,samples=70]
    (\x,{2.4*(1-(1-\x/5)*(1-\x/5))});
  \draw[blue!75,very thick] (5,2.4) -- (6.0,2.4);
  \fill[red!75] (5,2.4) circle (1.6pt);
  \fill[black!70] (6.0,2.4) circle (1.6pt);
  \node[above,font=\footnotesize,black!75,align=center] at (4.75,2.5) {最小充分表示};
  \node[right,font=\footnotesize,black!75] at (6.08,2.28) {\(T=X\)};
  \node[font=\footnotesize,black!60,align=center] at (3.4,2.78) {不可行};
  \node[font=\footnotesize,black!60,align=center] at (3.9,0.75) {可行区};
  \draw[black!45] (0.62,0.05) -- (0.62,0.53);
  \node[right,font=\footnotesize,black!60,align=left] at (0.72,0.32)
    {\(\beta\) 小：压得狠，\\任务信息也丢了};
\end{tikzpicture}

\emph{横轴是表示保留了多少输入信息，纵轴是它还能说明多少标签。可行区被两条线夹住：水平虚线是数据处理不等式给的天花板，斜虚线说明表示对标签的信息不可能超过它对输入的信息。蓝色前沿上每一点对应一个 \(\beta\)，切线斜率就是 \(1/\beta\)。红点右边那一整段水平线上的表示都是充分的，其中最靠左的一个最省——这就是最小充分表示。}
\end{center}

\subsection{哪些输入值该被合并}

前沿上的表示在做什么？一个特殊情形已经说明了大半。若两个输入值 \(x_1,x_2\) 的标签条件分布相同，

\[
p(y\given x_1)=p(y\given x_2),
\]

那么对预测 \(Y\) 而言它们不可区分。把它们映到同一个表示值，\(I(X;T)\) 下降，\(I(T;Y)\) 分毫不损——这是纯赚的压缩。

一般情形只是把``相同''放宽成``相近''。对 \(\mathcal L_{\mathrm{IB}}\) 在归一化约束下做变分，得到的平稳编码器满足自洽形式

\[
p(t\given x)\propto p(t)\exp\bigl[-\beta\,\KL(p(y\given x)\,\|\,p(y\given t))\bigr],
\]

它与 \(p(t)\) 和 \(p(y\given t)\) 的边缘式一起构成一组迭代更新。推导过程不必细看，看的是里面那个 KL：它量的是把 \(x\) 并入簇 \(t\) 之后，标签预测行为被扭曲了多少；扭曲小则并入的概率大。指数形式只是这句话的数学外壳。

于是信息瓶颈的解是一种软聚类，而聚类依据是输入对目标的预测行为，不是输入坐标之间的距离。两个像素分布截然不同的图片可以落进同一簇，只要它们对``猫还是狗''给出相同的证据。

\subsection{充分是理想边界，学到的只是近似}

把``一点不丢''写成条件独立：若 \(Y\perp X\given T\)，则 \(T\) 对预测 \(Y\) 充分，此时 \(I(T;Y)=I(X;Y)\)，表示保住了输入中全部任务相关的信息。在所有充分表示里再让 \(I(X;T)\) 最小，就是图上那个红点。同一个概念在统计里叫充分统计量，见\hyperref[sec:11-Mathematical-Statistics/03-Sufficient-Statistics-and-Information/01]{``充分统计量保留了哪些参数信息''}。

麻烦在于，前面 \(X=(Y,N)\) 的例子里 \(T=Y\) 之所以完美，是因为我们直接写下了 \(Y\)。真实预测场景里 \(Y\) 恰恰是要猜的东西，你只能从训练样本里学一个近似的编码器。问题于是从``联合分布的结构''变成``有限数据、有限容量、有限优化步数三者共同作用''，前沿的形状本身就不再确知。

\subsection{把它套到确定性网络上，会先撞到定义问题}

信息瓶颈作为优化原理是清楚的。可一旦把普通神经网络的隐藏层直接代进 \(I(X;T)\)，麻烦从定义层面就开始了，而且不是估计精度不够那种麻烦。

若 \(X\) 连续、\(T=f(X)\) 是确定性映射，联合分布会落在低维流形上，互信息可以取到无穷。若这个映射恰好可逆，理论互信息干脆恒等于 \(H(X)\)——网络怎么训练、几何形状看着变了多少，这个数一动不动。要让它变成有限且会动的量，你必须往里加东西：给表示加随机噪声，或者对输入和表示做量化。而加多少噪声、分多少个箱，直接决定了曲线长什么样。

这就把一个广为流传的经验主张放到了显微镜下。``深网训练分成先拟合、后压缩两个阶段''这条观察，在某些设置下能复现，在另一些设置下不能：换掉激活函数（双曲正切与 ReLU 的结论就不一致）、换掉分箱宽度、换掉互信息估计器，压缩阶段可能变弱甚至消失。也有工作论证观察到的``压缩''主要是饱和型激活函数的副产物，而非训练的普遍规律。

诚实的说法是：这是一个有争议的经验现象，不是已经定论的定律。要报告这类曲线，至少得同时交代隐藏表示有没有加随机噪声、输入与表示怎样量化、用了哪种估计器及其偏差方向，以及比较的双方是不是在同一尺度和同一随机变量定义下。这几项换一个，趋势就可能换一个方向。

\subsection{变分信息瓶颈：优化的是界，不是那个量}

既然精确互信息算不出来，实用做法是换成可计算的界，两侧各用一次同样的手法。

预测那一侧，\(I(T;Y)\) 里的 \(p(y\given t)\) 未知，就拿一个可训练的解码器 \(q(y\given t)\) 顶上。替换之后的差额恰好是一个 KL 散度，而 KL 非负，所以替换只会让读数变小——得到的是 \(I(T;Y)\) 的下界。压缩那一侧同理：\(I(X;T)\) 里的边缘 \(p(t)\) 拿一个固定的先验 \(r(t)\)（通常取标准 Gaussian）顶上，差额也是一个非负的 KL，于是 \(\E_X\KL(p(t\given X)\|r(t))\) 是 \(I(X;T)\) 的上界。

整套骨架就这一步：把算不出来的分布换成可训练的近似，用 Jensen 或者 KL 非负保证方向不出错。两个界合起来是一个能对参数求导的目标，形式上和\hyperref[sec:14-Deep-Learning/06-Variational-Autoencoders/01]{``VAE 把不可积后验变成可训练下界''}几乎一样，区别只在于 VAE 重建的是输入，这里预测的是标签。

限度要说清楚。训练时下降的是界，不是互信息本身，而界与真值之间的间隙由近似族的表达能力决定，通常无从测量。所以``变分信息瓶颈训练出来的模型压缩得更狠''这句话，严格说只意味着那个上界更小。把这个方法当成一种带信息论动机的正则化来用是稳的；把它的读数当成对真实互信息的测量则不稳。

信息瓶颈给出的是一个可写下来的表示目标，附带一条来自数据处理不等式的天花板：任何表示都不可能比 \(I(X;Y)\) 说得更多。可这条天花板还没有兑换成任何关于错误率的承诺——知道特征里只剩一点点标签信息，究竟意味着分类器最好能做到什么程度？下一节把这个上限翻译成错误率的下界。
